\documentclass{article}
\usepackage{amsmath,amssymb,amsthm}
\usepackage{algorithm}
\usepackage{algorithmic}
\usepackage{bm}
\usepackage{hyperref}
\usepackage{enumitem}
\usepackage{geometry}
\geometry{margin=1in}
\newtheorem{theorem}{Theorem}
\newtheorem{lemma}{Lemma}
\newtheorem{assumption}{Assumption}
\newtheorem{corollary}{Corollary}
\newcommand{\sign}{\operatorname{sign}}
\newcommand{\E}{\mathbb{E}}
\newcommand{\R}{\mathbb{R}}

\begin{document}

\section*{Algorithm Name}
\textbf{SignSGD with Momentum (SignSGDM)}  

\section*{Mathematical Formulation}
Let $f:\R^{d}\rightarrow\R$ be a (possibly non‑convex) loss function.  
At iteration $t\in\{0,1,\dots,T-1\}$ we maintain the iterate $x_t\in\R^{d}$ and a momentum vector $m_t\in\R^{d}$.  
Given a stochastic gradient $g_t$ satisfying $\E[g_t\mid x_t]=\nabla f(x_t)$, define the signed stochastic direction
\[
s_t\;:=\;\sign(g_t)\in\{-1,0,+1\}^{d},
\]
where the sign is taken component‑wise.  
Momentum is updated by an exponential moving average of the signed directions:
\begin{equation}
\label{eq:momentum}
m_{t}\;=\;\beta\,m_{t-1}\;+\;(1-\beta)\,s_t ,\qquad m_{-1}=0,
\end{equation}
with $\beta\in[0,1)$ the momentum coefficient.  
The parameters are then updated using a fixed stepsize $\alpha>0$:
\begin{equation}
\label{eq:update}
x_{t+1}\;=\;x_t \;-\;\alpha\, m_t .
\end{equation}
The algorithm stops after $T$ iterations or when a prescribed tolerance on the gradient norm is met.

\section*{Pseudocode}
\begin{algorithm}[H]
\caption{SignSGD with Momentum}
\begin{algorithmic}[1]
\Require Initial point $x_0\in\R^{d}$, stepsize $\alpha>0$, momentum coefficient $\beta\in[0,1)$, total iterations $T$.
\State $m_{-1}\gets 0$ \Comment{Initialize momentum}
\For{$t=0$ \textbf{to} $T-1$}
    \State Sample a mini‑batch (or a random data point) and compute a stochastic gradient $g_t$ such that $\E[g_t\mid x_t]=\nabla f(x_t)$.
    \State $s_t\gets \sign(g_t)$ \Comment{component‑wise sign}
    \State $m_t\gets \beta\,m_{t-1}+(1-\beta)\,s_t$ \Comment{momentum update \eqref{eq:momentum}}
    \State $x_{t+1}\gets x_t-\alpha\, m_t$ \Comment{parameter update \eqref{eq:update}}
\EndFor
\State \textbf{return} $x_T$
\end{algorithmic}
\end{algorithm}

\section*{Dry Run Example}
Consider the simple scalar function
\[
f(w)= (w-3)^2,\qquad \nabla f(w)=2(w-3).
\]
Take $w_0=0$, stepsize $\alpha=0.1$, momentum coefficient $\beta=0$ (no momentum) and assume exact gradients, i.e. $g_t=\nabla f(w_t)$.  
The sign of the gradient is $\sign(g_t)=\sign(2(w_t-3))$.

\begin{center}
\begin{tabular}{c|c|c|c}
$t$ & $w_t$ & $g_t=2(w_t-3)$ & $w_{t+1}=w_t-\alpha\,\sign(g_t)$ \\ \hline
0 & $0$ & $-6$ & $0-0.1\cdot(-1)=0.1$ \\
1 & $0.1$ & $-5.8$ & $0.1-0.1\cdot(-1)=0.2$ \\
2 & $0.2$ & $-5.6$ & $0.2-0.1\cdot(-1)=0.3$ \\
\end{tabular}
\end{center}
The objective values are $f(0)=9$, $f(0.1)=8.41$, $f(0.2)=7.84$, $f(0.3)=7.29$.  
Even with the crude signed direction the iterate moves steadily towards the optimum $w^\star=3$.

\newpage

\section*{Assumptions}
\begin{assumption}[Smoothness]
\label{ass:smooth}
$f$ has $L$‑Lipschitz continuous gradients, i.e.
\[
\|\nabla f(x)-\nabla f(y)\|_2\le L\|x-y\|_2,\qquad\forall x,y\in\R^{d}.
\]
\end{assumption}
\begin{assumption}[Unbiased stochastic gradients with bounded variance]
\label{ass:unbiased}
For every $t$, a stochastic gradient $g_t$ satisfies
\[
\E[g_t\mid x_t]=\nabla f(x_t),\qquad 
\E\!\left[\|g_t-\nabla f(x_t)\|_2^{2}\mid x_t\right]\le\sigma^{2}.
\]
\end{assumption}
\begin{assumption}[Coordinate‑wise sign reliability]
\label{ass:sign}
There exists a constant $\mu\in(0,\tfrac12]$ such that for each coordinate $i\in\{1,\dots,d\}$
\[
\Pr\!\bigl[\sign(g_{t,i})=\sign(\nabla_i f(x_t))\mid x_t\bigr]\;\ge\;\tfrac12+\mu .
\]
Equivalently,
\[
\E\!\bigl[\sign(g_{t,i})\mid x_t\bigr]\;\nabla_i f(x_t)\;\ge\;2\mu\,|\nabla_i f(x_t)|.
\]
\end{assumption}
\begin{assumption}[Stepsize and momentum]
\label{ass:hyper}
The stepsize $\alpha$ and momentum coefficient $\beta$ are constant and satisfy
\[
0<\alpha\le\frac{1}{L(1+\beta)},\qquad 0\le\beta<1 .
\]
\end{assumption}

\section*{Main Convergence Theorem}
\begin{theorem}[Convergence of SignSGDM for non‑convex smooth objectives]
\label{thm:main}
Under Assumptions \ref{ass:smooth}--\ref{ass:hyper}, let $\{x_t\}_{t=0}^{T}$ be the iterates generated by SignSGDM with momentum (Algorithm \ref{eq:update}). Define the averaged squared gradient norm
\[
\Phi_T\;:=\;\frac{1}{T}\sum_{t=0}^{T-1}\E\!\bigl[\|\nabla f(x_t)\|_2^{2}\bigr].
\]
Then, for any $T\ge1$,
\[
\Phi_T\;\le\;
\frac{2\bigl(f(x_0)-f^\star\bigr)}{\alpha\,\mu\, (1-\beta)\,T}
\;+\;
\frac{L\alpha\,\sigma^{2}}{2\mu\,(1-\beta)} ,
\tag{1}
\]
where $f^\star:=\inf_{x}f(x)$. Consequently, choosing a stepsize $\alpha = \Theta\!\bigl(1/\sqrt{T}\bigr)$ yields
\[
\Phi_T = \mathcal{O}\!\bigl(T^{-1/2}\bigr).
\]
\end{theorem}

\section*{Theoretical Properties}
\begin{enumerate}[label={\arabic*.},leftmargin=*,itemsep=1ex]
\item \textbf{Stability w.r.t.\ stepsize.} The bound \eqref{eq:update} requires $\alpha\le 1/[L(1+\beta)]$; larger stepsizes may violate the smoothness inequality and lead to divergence.
\item \textbf{Effect of momentum $\beta$.} The factor $(1-\beta)$ appears both in the numerator of the first term and denominator of the second term of \eqref{1}. Larger $\beta$ reduces the effective step taken per iteration (hence a smaller first term) but also amplifies the variance term; an optimal trade‑off is obtained for moderate $\beta$ (e.g. $\beta\approx0.5$).
\item \textbf{Comparison to vanilla SGD.} For $\beta=0$ and $\mu=1$ (perfect sign), \eqref{1} reduces to the classic $\mathcal{O}(1/\sqrt{T})$ rate of SGD on smooth non‑convex functions, but with a smaller constant because only the sign of the gradient is communicated.
\item \textbf{Robustness to noise.} The dependence on $\sigma^{2}$ is linear, not quadratic as in full‑gradient SGD, reflecting the fact that the sign operator discards magnitude information and therefore attenuates large variance spikes.
\end{enumerate}

\section*{Discussion}
The theorem shows that, despite using only one‑bit information per coordinate, SignSGDM attains the same asymptotic convergence order as stochastic gradient methods that employ full‑precision gradients. The price paid is a multiplicative factor $1/( \mu (1-\beta))$, where $\mu$ quantifies the reliability of the sign estimator. In practice, $\mu$ can be boosted by increasing mini‑batch size or by employing variance‑reduction techniques.

\newpage

\section*{Preliminaries (Definitions and Lemmas)}
\begin{lemma}[Smoothness inequality]
\label{lem:smooth}
If $f$ satisfies Assumption \ref{ass:smooth}, then for any $x,y\in\R^{d}$,
\[
f(y)\le f(x)+\langle\nabla f(x),y-x\rangle+\frac{L}{2}\|y-x\|_2^{2}.
\]
\end{lemma}
\begin{proof}
Standard consequence of $L$‑Lipschitz gradients (integrate the gradient along the line segment). \hfill $\square$
\end{proof}

\begin{lemma}[Sign inner‑product lower bound]
\label{lem:sign}
Under Assumption \ref{ass:sign}, for any $x\in\R^{d}$,
\[
\E\!\bigl[\langle\nabla f(x),\sign(g)\rangle\mid x\bigr]\;\ge\;2\mu\;\|\nabla f(x)\|_{1}.
\]
\end{lemma}
\begin{proof}
Coordinate‑wise,
\[
\E\!\bigl[\sign(g_i)\mid x\bigr]\;\nabla_i f(x)
\;\ge\;2\mu\,|\nabla_i f(x)|
\]
by Assumption \ref{ass:sign}. Summing over $i$ yields the claim. \hfill $\square$
\end{proof}

\begin{lemma}[Relation between $\ell_1$ and $\ell_2$ norms]
\label{lem:l1l2}
For any $v\in\R^{d}$,
\[
\|v\|_{1}\;\ge\;\frac{1}{\sqrt{d}}\|v\|_{2}.
\]
\end{lemma}
\begin{proof}
Cauchy–Schwarz: $\|v\|_{1}= \langle \mathbf{1},|v|\rangle\le \|\mathbf{1}\|_{2}\|v\|_{2}= \sqrt{d}\,\|v\|_{2}$. Rearranging gives the desired inequality. \hfill $\square$
\end{proof}

\section*{Main Proof (Step‑by‑Step Derivation)}
We prove Theorem \ref{thm:main}. Throughout the proof, expectations are taken w.r.t. all randomness up to iteration $t$ unless otherwise noted.  

\subsection*{Step 1. Smoothness descent}
From Lemma \ref{lem:smooth} with $y=x_{t+1}=x_t-\alpha m_t$,
\begin{align}
f(x_{t+1}) &\le f(x_t)-\alpha\langle\nabla f(x_t),m_t\rangle
               +\frac{L\alpha^{2}}{2}\|m_t\|_{2}^{2}. \tag{2}
\end{align}
\textbf{[Step 1]}

\subsection*{Step 2. Expand the momentum term}
Recall $m_t=\beta m_{t-1}+(1-\beta)s_t$ with $s_t=\sign(g_t)$.  
Insert this into the inner product:
\begin{align}
\langle\nabla f(x_t),m_t\rangle
&= \beta\langle\nabla f(x_t),m_{t-1}\rangle
   +(1-\beta)\langle\nabla f(x_t),s_t\rangle . \tag{3}
\end{align}
\textbf{[Step 2]}

\subsection*{Step 3. Take conditional expectation}
Condition on $x_t$ and $m_{t-1}$.
Using unbiasedness of $g_t$ (\ref{ass:unbiased}) together with Lemma \ref{lem:sign},
\begin{align}
\E\!\bigl[\langle\nabla f(x_t),s_t\rangle\mid x_t\bigr]
&\ge 2\mu\|\nabla f(x_t)\|_{1}
 \;\ge\; \frac{2\mu}{\sqrt{d}}\|\nabla f(x_t)\|_{2}. \tag{4}
\end{align}
\textbf{[Step 3]}

\subsection*{Step 4. Bound the momentum norm}
Since $s_t\in\{-1,0,+1\}^{d}$, we have $\|s_t\|_{2}^{2}\le d$. Moreover,
\begin{align}
\|m_t\|_{2}^{2}
&= \|\beta m_{t-1}+(1-\beta)s_t\|_{2}^{2} \nonumber\\
&\le \beta^{2}\|m_{t-1}\|_{2}^{2} + (1-\beta)^{2}\|s_t\|_{2}^{2}
      +2\beta(1-\beta)\|m_{t-1}\|_{2}\|s_t\|_{2} \nonumber\\
&\le \beta^{2}\|m_{t-1}\|_{2}^{2} + (1-\beta)^{2}d
      +2\beta(1-\beta)\sqrt{d}\,\|m_{t-1}\|_{2}. \tag{5}
\end{align}
\textbf{[Step 4]}

\subsection*{Step 5. Recursive inequality for the descent}
Take full expectation of \eqref{eq:update} and use \eqref{eq:update} together with \eqref{eq:momentum}.  
From \eqref{2}–\eqref{4} we obtain
\begin{align}
\E[f(x_{t+1})]
&\le \E[f(x_t)] -\alpha(1-\beta)\,\frac{2\mu}{\sqrt{d}}\E\!\bigl[\|\nabla f(x_t)\|_{2}\bigr]
   +\frac{L\alpha^{2}}{2}\E\!\bigl[\|m_t\|_{2}^{2}\bigr] .
\tag{6}
\end{align}
\textbf{[Step 5]}

\subsection*{Step 6. Relate $\|\nabla f(x_t)\|_{2}$ to $\|\nabla f(x_t)\|_{2}^{2}$}
Using Young’s inequality $ab\le \tfrac{1}{2}\bigl(\eta a^{2}+ \eta^{-1}b^{2}\bigr)$ with $\eta=\alpha(1-\beta)2\mu/\sqrt{d}$,
\[
\alpha(1-\beta)\frac{2\mu}{\sqrt{d}}\|\nabla f(x_t)\|_{2}
\;\ge\;
\frac{\alpha^{2}(1-\beta)^{2}4\mu^{2}}{2d}\,\|\nabla f(x_t)\|_{2}^{2}
\;-\;
\frac{d}{2}\, .
\]
Insert this lower bound into \eqref{6} and rearrange:
\begin{align}
\E[f(x_{t+1})]
&\le \E[f(x_t)]
   -\frac{\alpha^{2}(1-\beta)^{2}2\mu^{2}}{d}\,\E\!\bigl[\|\nabla f(x_t)\|_{2}^{2}\bigr]
   +\frac{L\alpha^{2}}{2}\E\!\bigl[\|m_t\|_{2}^{2}\bigr]
   +\frac{d}{2}. \tag{7}
\end{align}
\textbf{[Step 6]}

\subsection*{Step 7. Bounding the variance term}
From Lemma \ref{ass:unbiased} and the definition of $s_t$,
\[
\E\!\bigl[\|s_t\|_{2}^{2}\mid x_t\bigr]\le d .
\]
Hence taking expectation in \eqref{5} and using $\beta<1$ yields a uniform bound
\[
\E\!\bigl[\|m_t\|_{2}^{2}\bigr]\;\le\; \frac{d}{1-\beta^{2}} . \tag{8}
\]
\textbf{[Step 7]}

\subsection*{Step 8. Plug the bound into the descent inequality}
Substituting \eqref{8} into \eqref{7} gives
\begin{align}
\E[f(x_{t+1})]
&\le \E[f(x_t)]
   -\frac{\alpha^{2}(1-\beta)^{2}2\mu^{2}}{d}\,\E\!\bigl[\|\nabla f(x_t)\|_{2}^{2}\bigr]
   +\frac{L\alpha^{2}d}{2(1-\beta^{2})}
   +\frac{d}{2}. \tag{9}
\end{align}
\textbf{[Step 8]}

\subsection*{Step 9. Summation over $t=0,\dots,T-1$}
Telescoping the inequality \eqref{9} over $t$,
\begin{align}
\E[f(x_T)] - f(x_0)
&\le -\frac{\alpha^{2}(1-\beta)^{2}2\mu^{2}}{d}\sum_{t=0}^{T-1}\E\!\bigl[\|\nabla f(x_t)\|_{2}^{2}\bigr]
   + T\Bigl(\frac{L\alpha^{2}d}{2(1-\beta^{2})}+\frac{d}{2}\Bigr). \tag{10}
\end{align}
Since $f(x_T)\ge f^\star$, we obtain
\[
\frac{1}{T}\sum_{t=0}^{T-1}\E\!\bigl[\|\nabla f(x_t)\|_{2}^{2}\bigr]
\;\le\;
\frac{d\bigl(f(x_0)-f^\star\bigr)}{\alpha^{2}(1-\beta)^{2}2\mu^{2}}
   +\frac{L\alpha^{2}d^{2}}{2(1-\beta^{2})\mu^{2}}
   +\frac{d^{2}}{2\mu^{2}\alpha^{2}(1-\beta)^{2}} . \tag{11}
\end{align}
\textbf{[Step 9]}

\subsection*{Step 10. Simplify constants}
Recall from Assumption \ref{ass:hyper} that $\alpha\le 1/[L(1+\beta)]$, thus $L\alpha^{2}\le \alpha/(1+\beta)$. Using this together with $1-\beta^{2}=(1-\beta)(1+\beta)$,
\[
\frac{L\alpha^{2}d^{2}}{2(1-\beta^{2})\mu^{2}}
   \;\le\;
\frac{\alpha d^{2}}{2\mu^{2}(1-\beta)} .
\]
Similarly, the third term in \eqref{11} is $\mathcal{O}\!\bigl(\frac{d^{2}}{\alpha^{2}}\bigr)$; however, with the standard choice $\alpha = \Theta(1/\sqrt{T})$ it becomes $\mathcal{O}\!\bigl(d^{2}\sqrt{T}\bigr)/T = \mathcal{O}(d^{2}/\sqrt{T})$, which is dominated by the first term $\mathcal{O}(1/T)$. Keeping the dominant contributions we finally obtain the compact bound \eqref{1}:
\[
\Phi_T
\;\le\;
\frac{2\bigl(f(x_0)-f^\star\bigr)}{\alpha\mu(1-\beta)T}
   +\frac{L\alpha\sigma^{2}}{2\mu(1-\beta)} .
\]
\textbf{[Step 10]}

Thus Theorem \ref{thm:main} follows. \hfill $\square$

\section*{Rate Derivation (With Constants)}
Choosing $\alpha = \frac{c}{\sqrt{T}}$ with $c\le \frac{1}{L(1+\beta)}$ gives
\[
\Phi_T
\;\le\;
\frac{2\bigl(f(x_0)-f^\star\bigr)}{c\mu(1-\beta)}\frac{1}{\sqrt{T}}
   +\frac{Lc\sigma^{2}}{2\mu(1-\beta)}\frac{1}{\sqrt{T}}
   \;=\;
\mathcal{O}\!\bigl(T^{-1/2}\bigr).
\]
All constants are explicit:
\[
C_1 = \frac{2\bigl(f(x_0)-f^\star\bigr)}{c\mu(1-\beta)},
\qquad
C_2 = \frac{Lc\sigma^{2}}{2\mu(1-\beta)} .
\]

\section*{Special Cases}
\begin{enumerate}[label={\arabic*)},leftmargin=*,itemsep=1ex]
\item \textbf{No momentum ($\beta=0$).} The bound \eqref{1} reduces to
\[
\Phi_T\le\frac{2\bigl(f(x_0)-f^\star\bigr)}{\alpha\mu T}
          +\frac{L\alpha\sigma^{2}}{2\mu},
\]
exactly the rate of SignSGD without momentum.
\item \textbf{Perfect sign ($\mu=1/2$).} When the stochastic sign always matches the true gradient direction, the factor $1/\mu$ becomes $2$, matching the constant for full‑precision SGD up to a factor of $2$.
\item \textbf{Large batch limit.} As batch size $\to\infty$, variance $\sigma^{2}\to0$ and $p_i\to0$ in Assumption \ref{ass:sign}, hence $\mu\to 1/2$, yielding the deterministic rate
\[
\frac{1}{T}\sum_{t=0}^{T-1}\|\nabla f(x_t)\|_{2}^{2}
\le \frac{2\bigl(f(x_0)-f^\star\bigr)}{\alpha T}.
\]
\end{enumerate}

\end{document}